% Worked Examples: Topic 2.5 - Transfer Functions
\documentclass[11pt,a4paper]{article}
\usepackage{../worked-examples-template}

\title{\textbf{Worked Examples}\\[0.5em]
\large Topic 2.5: Transfer Functions\\[0.3em]
\normalsize \coursesubtitle{} -- \coursetitle}
\author{Assoc.\ Prof.\ William Robertson \and Dr Sean McGowan}
\date{\today}

\begin{document}

\maketitle
\tableofcontents
\newpage

\section{Concept 2.5.1: Frequency Domain}

\subsection{Example 1: Frequency Response of a First-Order System}

\paragraph{Problem:} Consider a first-order system with transfer function:
\begin{align}
G(s) = \frac{5}{s+2}
\end{align}

Find the steady-state response to a sinusoidal input $u(t) = 3\sin(4t)$ (amplitude 3, frequency $\omega = 4$ rad/s).

\paragraph{Solution:}

For a sinusoidal input $u(t) = A\sin(\omega t)$, the steady-state output is:
\begin{align}
y_{ss}(t) = A|G(j\omega)|\sin(\omega t + \angle G(j\omega))
\end{align}

where $|G(j\omega)|$ is the magnitude and $\angle G(j\omega)$ is the phase of the transfer function at frequency $\omega$.

For $\omega = 4$ rad/s:
\begin{align}
G(j4) = \frac{5}{j4+2} = \frac{5}{2+j4}
\end{align}

To find magnitude and phase, convert to polar form:
\begin{align}
G(j4) &= \frac{5}{2+j4} \cdot \frac{2-j4}{2-j4} = \frac{5(2-j4)}{4+16} = \frac{10-j20}{20} = 0.5 - j
\end{align}

Magnitude:
\begin{align}
|G(j4)| = \sqrt{(0.5)^2 + (-1)^2} = \sqrt{0.25 + 1} = \sqrt{1.25} \approx 1.118
\end{align}

Phase:
\begin{align}
\angle G(j4) = \arctan\left(\frac{-1}{0.5}\right) = \arctan(-2) \approx -63.43° \approx -1.107 \text{ rad}
\end{align}

Therefore, the steady-state output is:
\begin{align}
y_{ss}(t) &= 3 \times 1.118 \times \sin(4t - 1.107) \\
&\approx 3.35\sin(4t - 63.43°)
\end{align}

\textbf{Interpretation:} The output amplitude is attenuated by a factor of 1.118 (from 3 to 3.35), and there is a phase lag of approximately 63.43° (or 1.107 radians).

\section{Concept 2.5.2: Transfer Functions}

\subsection{Example 1: Deriving Transfer Function from State-Space Model}

\paragraph{Problem:} Find the transfer function $G(s) = Y(s)/U(s)$ for the state-space system:
\begin{align}
\frac{dx}{dt} = \begin{bmatrix} -2 & 1 \\ 0 & -3 \end{bmatrix}x + \begin{bmatrix} 1 \\ 1 \end{bmatrix}u, \quad y = \begin{bmatrix} 1 & 2 \end{bmatrix}x
\end{align}

\paragraph{Solution:}

The transfer function from state-space is given by:
\begin{align}
G(s) = C(sI - A)^{-1}B + D
\end{align}

where $D = 0$ (no direct feedthrough).

First, compute $sI - A$:
\begin{align}
sI - A = \begin{bmatrix} s & 0 \\ 0 & s \end{bmatrix} - \begin{bmatrix} -2 & 1 \\ 0 & -3 \end{bmatrix} = \begin{bmatrix} s+2 & -1 \\ 0 & s+3 \end{bmatrix}
\end{align}

Find the inverse using the adjugate method. For a $2 \times 2$ matrix, the adjugate (also called adjoint) is the transpose of the cofactor matrix:
\begin{align}
(sI-A)^{-1} = \frac{1}{\det(sI-A)}\text{adj}(sI-A)
\end{align}

Determinant:
\begin{align}
\det(sI-A) = (s+2)(s+3) = s^2 + 5s + 6
\end{align}

Adjugate (swap diagonal, negate off-diagonal):
\begin{align}
\text{adj}(sI-A) = \begin{bmatrix} s+3 & 1 \\ 0 & s+2 \end{bmatrix}
\end{align}

Therefore:
\begin{align}
(sI-A)^{-1} = \frac{1}{s^2+5s+6}\begin{bmatrix} s+3 & 1 \\ 0 & s+2 \end{bmatrix}
\end{align}

Compute $(sI-A)^{-1}B$:
\begin{align}
(sI-A)^{-1}B &= \frac{1}{s^2+5s+6}\begin{bmatrix} s+3 & 1 \\ 0 & s+2 \end{bmatrix}\begin{bmatrix} 1 \\ 1 \end{bmatrix} \\
&= \frac{1}{s^2+5s+6}\begin{bmatrix} s+4 \\ s+2 \end{bmatrix}
\end{align}

Finally, multiply by $C$:
\begin{align}
G(s) &= C(sI-A)^{-1}B \\
&= \begin{bmatrix} 1 & 2 \end{bmatrix}\frac{1}{s^2+5s+6}\begin{bmatrix} s+4 \\ s+2 \end{bmatrix} \\
&= \frac{(s+4) + 2(s+2)}{s^2+5s+6} \\
&= \frac{s+4+2s+4}{s^2+5s+6} \\
&= \frac{3s+8}{s^2+5s+6} \\
&= \frac{3s+8}{(s+2)(s+3)}
\end{align}

\textbf{Result:} The transfer function is:
\begin{align}
G(s) = \frac{3s+8}{(s+2)(s+3)}
\end{align}

The poles are at $s = -2$ and $s = -3$ (eigenvalues of $A$).
The zero is at $s = -8/3 \approx -2.67$.

\section{Concept 2.5.3: Laplace Transform}

\subsection{Example 1: Solving an ODE using Laplace Transforms}

\paragraph{Problem:} Use Laplace transforms to solve the differential equation:
\begin{align}
\ddot{y} + 5\dot{y} + 6y = 2u(t)
\end{align}

where $u(t)$ is a unit step function and initial conditions are $y(0) = 0$, $\dot{y}(0) = 0$.

\paragraph{Solution:}

Taking the Laplace transform of both sides (assuming zero initial conditions):
\begin{align}
s^2Y(s) + 5sY(s) + 6Y(s) = \frac{2}{s}
\end{align}

Note: $\mathcal{L}\{u(t)\} = \mathcal{L}\{1\} = \frac{1}{s}$ for a unit step.

Factor out $Y(s)$:
\begin{align}
Y(s)(s^2 + 5s + 6) = \frac{2}{s}
\end{align}

Solve for $Y(s)$:
\begin{align}
Y(s) = \frac{2}{s(s^2+5s+6)} = \frac{2}{s(s+2)(s+3)}
\end{align}

Use partial fraction decomposition:
\begin{align}
\frac{2}{s(s+2)(s+3)} = \frac{A}{s} + \frac{B}{s+2} + \frac{C}{s+3}
\end{align}

Multiply both sides by $s(s+2)(s+3)$:
\begin{align}
2 = A(s+2)(s+3) + Bs(s+3) + Cs(s+2)
\end{align}

Find coefficients:
\begin{itemize}
\item Set $s = 0$: $2 = A(2)(3) = 6A \Rightarrow A = \frac{1}{3}$
\item Set $s = -2$: $2 = B(-2)(-2+3) = -2B \Rightarrow B = -1$
\item Set $s = -3$: $2 = C(-3)(-3+2) = 3C \Rightarrow C = \frac{2}{3}$
\end{itemize}

Therefore:
\begin{align}
Y(s) = \frac{1/3}{s} + \frac{-1}{s+2} + \frac{2/3}{s+3}
\end{align}

Take the inverse Laplace transform using the table:
\begin{itemize}
\item $\mathcal{L}^{-1}\left\{\frac{1}{s}\right\} = 1$
\item $\mathcal{L}^{-1}\left\{\frac{1}{s+a}\right\} = e^{-at}$
\end{itemize}

Result:
\begin{align}
y(t) = \frac{1}{3} - e^{-2t} + \frac{2}{3}e^{-3t}, \quad t \geq 0
\end{align}

\textbf{Verification:} As $t \to \infty$, the exponential terms decay to zero:
\begin{align}
\lim_{t \to \infty} y(t) = \frac{1}{3}
\end{align}

This matches the expected steady-state value from the final value theorem:
\begin{align}
y_{ss} = \lim_{s \to 0} sY(s) = \lim_{s \to 0} s \cdot \frac{2}{s(s+2)(s+3)} = \frac{2}{2 \cdot 3} = \frac{1}{3}
\end{align}

\section{Concept 2.5.4: Block Diagrams}

\subsection{Example 1: Block Diagram Reduction}

\paragraph{Problem:} Consider a feedback control system with the following components in series and feedback configuration:
\begin{itemize}
\item Controller: $C(s) = K_p + \frac{K_i}{s}$ (PI controller)
\item Plant: $P(s) = \frac{1}{s+1}$
\item Feedback sensor: $H(s) = 1$ (unity feedback)
\end{itemize}

Find the closed-loop transfer function from reference $R(s)$ to output $Y(s)$.

\paragraph{Solution:}

For a standard feedback configuration with controller $C(s)$, plant $P(s)$, and feedback $H(s) = 1$, the closed-loop transfer function is:
\begin{align}
T(s) = \frac{C(s)P(s)}{1 + C(s)P(s)H(s)} = \frac{C(s)P(s)}{1 + C(s)P(s)}
\end{align}

First, find the product $C(s)P(s)$:
\begin{align}
C(s)P(s) &= \left(K_p + \frac{K_i}{s}\right) \cdot \frac{1}{s+1} \\
&= \frac{K_ps + K_i}{s} \cdot \frac{1}{s+1} \\
&= \frac{K_ps + K_i}{s(s+1)}
\end{align}

Now compute the closed-loop transfer function:
\begin{align}
T(s) &= \frac{\frac{K_ps + K_i}{s(s+1)}}{1 + \frac{K_ps + K_i}{s(s+1)}} \\
&= \frac{K_ps + K_i}{s(s+1) + K_ps + K_i} \\
&= \frac{K_ps + K_i}{s^2 + s + K_ps + K_i} \\
&= \frac{K_ps + K_i}{s^2 + (1+K_p)s + K_i}
\end{align}

\textbf{Result:} The closed-loop transfer function is:
\begin{align}
T(s) = \frac{K_ps + K_i}{s^2 + (1+K_p)s + K_i}
\end{align}

\textbf{Analysis:}
\begin{itemize}
\item The system is second-order with one zero at $s = -K_i/K_p$
\item The closed-loop poles depend on values of $K_p$ and $K_i$
\item For a specific example with $K_p = 2$ and $K_i = 3$:
\begin{align}
T(s) = \frac{2s + 3}{s^2 + 3s + 3}
\end{align}
The poles are at $s = \frac{-3 \pm \sqrt{9-12}}{2} = \frac{-3 \pm j\sqrt{3}}{2} \approx -1.5 \pm j0.866$
\end{itemize}

\section{Concept 2.5.5: DC Gain}

\subsection{Example 1: Computing DC Gain and Steady-State Response}

\paragraph{Problem:} For the transfer function:
\begin{align}
G(s) = \frac{10(s+2)}{(s+1)(s+5)}
\end{align}

Find: (a) the DC gain, and (b) the steady-state output for a unit step input.

\paragraph{Solution:}

\textbf{Part (a): DC Gain}

The DC gain is the value of the transfer function at $s = 0$:
\begin{align}
G(0) = \frac{10(0+2)}{(0+1)(0+5)} = \frac{20}{5} = 4
\end{align}

\textbf{Part (b): Steady-State Response to Unit Step}

For a unit step input, $U(s) = \frac{1}{s}$, the output is:
\begin{align}
Y(s) = G(s)U(s) = \frac{10(s+2)}{(s+1)(s+5)} \cdot \frac{1}{s}
\end{align}

Using the final value theorem:
\begin{align}
y_{ss} = \lim_{t \to \infty} y(t) = \lim_{s \to 0} sY(s)
\end{align}

Compute:
\begin{align}
y_{ss} &= \lim_{s \to 0} s \cdot \frac{10(s+2)}{s(s+1)(s+5)} \\
&= \lim_{s \to 0} \frac{10(s+2)}{(s+1)(s+5)} \\
&= \frac{10(2)}{(1)(5)} = \frac{20}{5} = 4
\end{align}

\textbf{Result:} 
\begin{itemize}
\item DC gain: $G(0) = 4$
\item Steady-state step response: $y_{ss} = 4$
\end{itemize}

\textbf{Interpretation:} The DC gain equals the steady-state response to a unit step input, which is a fundamental property of stable systems. An input step of magnitude 1 results in an output that settles to a value of 4.

\section{Concept 2.5.6: Bode Plots}

\subsection{Example 1: Sketching Bode Plots by Hand}

\paragraph{Problem:} Sketch the Bode magnitude and phase plots for the transfer function:
\begin{align}
G(s) = \frac{100}{s(s+10)}
\end{align}

\paragraph{Solution:}

First, rewrite in standard Bode form by factoring out constants and normalizing poles/zeros:
\begin{align}
G(s) = \frac{100}{s \cdot 10(s/10 + 1)} = \frac{10}{s(1 + s/10)}
\end{align}

At $s = j\omega$:
\begin{align}
G(j\omega) = \frac{10}{j\omega(1 + j\omega/10)}
\end{align}

\textbf{Magnitude Plot:}

Express in dB:
\begin{align}
|G(j\omega)|_{dB} = 20\log_{10}|G(j\omega)| = 20\log_{10}\left|\frac{10}{j\omega(1 + j\omega/10)}\right|
\end{align}

Using logarithm properties:
\begin{align}
|G(j\omega)|_{dB} &= 20\log_{10}(10) - 20\log_{10}|\omega| - 20\log_{10}|1 + j\omega/10| \\
&= 20 - 20\log_{10}(\omega) - 20\log_{10}\sqrt{1 + (\omega/10)^2}
\end{align}

Asymptotic approximation:
\begin{itemize}
\item \textbf{Integrator $1/s$:} Contributes $-20$ dB/decade slope starting from DC
\item \textbf{Pole at $\omega = 10$ rad/s:} Adds $-20$ dB/decade for $\omega > 10$
\item \textbf{DC gain term (10):} Contributes $+20$ dB offset
\end{itemize}

Key frequencies:
\begin{itemize}
\item At $\omega = 1$ rad/s: $|G(j1)|_{dB} = 20 - 0 - 0 = 20$ dB
\item At $\omega = 10$ rad/s (corner frequency): $|G(j10)|_{dB} \approx 20 - 20 - 3 = -3$ dB
\item At $\omega = 100$ rad/s: $|G(j100)|_{dB} = 20 - 40 - 20 = -40$ dB
\end{itemize}

Slope breakdown:
\begin{itemize}
\item $\omega < 10$: slope is $-20$ dB/decade (from integrator)
\item $\omega > 10$: slope is $-40$ dB/decade (integrator + pole)
\end{itemize}

\textbf{Phase Plot:}

Phase contribution from each component:
\begin{itemize}
\item Integrator $1/s$: contributes $-90°$ at all frequencies
\item Pole at $s = -10$: contributes $-\arctan(\omega/10)$
  \begin{itemize}
  \item $\omega \ll 10$: $\approx 0°$
  \item $\omega = 10$: $\approx -45°$
  \item $\omega \gg 10$: $\approx -90°$
  \end{itemize}
\end{itemize}

Total phase:
\begin{align}
\angle G(j\omega) = -90° - \arctan(\omega/10)
\end{align}

Key phase values:
\begin{itemize}
\item At $\omega = 1$: $\angle G \approx -90° - 6° = -96°$
\item At $\omega = 10$: $\angle G = -90° - 45° = -135°$
\item At $\omega = 100$: $\angle G \approx -90° - 84° = -174°$
\item As $\omega \to \infty$: $\angle G \to -180°$
\end{itemize}

\textbf{Summary:}
\begin{itemize}
\item Magnitude starts at 20 dB at $\omega = 1$ rad/s with $-20$ dB/decade slope
\item At corner frequency $\omega = 10$ rad/s, slope changes to $-40$ dB/decade
\item Phase starts at $-90°$ and asymptotically approaches $-180°$
\item Phase is $-135°$ at the corner frequency
\end{itemize}

\end{document}
